
\subsection{\S~312.45 Inactive Status}

(a) If no subjects are entered into clinical studies for a period of 2 years or more under an IND, or if all investigations under an IND remain on clinical hold for 1 year or more, the IND may be placed by FDA on inactive status. This action may be taken by FDA either on request of the sponsor or on FDA's own initiative. If FDA seeks to act on its own initiative, it shall first notify the sponsor in writing, and the sponsor shall have 30 days to respond as to why the IND should continue to remain active.

(b) If an IND is placed on inactive status, all investigators shall be so notified and all stocks of the drug shall be returned or otherwise disposed of in accordance with \S~312.59.

(c) A sponsor is not required to submit annual reports to an IND on inactive status. An inactive IND is, however, still in effect for purposes of the public disclosure of data and information under \S~312.130.

(d) A sponsor who intends to resume clinical investigation under an IND placed on inactive status shall submit a protocol amendment under \S~312.30 containing the proposed general investigational plan and appropriate protocols. Clinical investigations under an IND on inactive status may only resume 30 days after FDA receives the protocol amendment, unless FDA notifies the sponsor that the investigations are subject to a clinical hold, or on earlier notification by FDA.

(e) An IND that remains on inactive status for 5 years or more may be terminated under \S~312.44.

\subsubsection{Physical AI Adaptation of \S~312.45}

(f) \textit{Physical AI system dormancy.} When an IND is placed on inactive status and the investigation involved Physical AI systems, the sponsor shall:

\begin{description}[leftmargin=0.8cm,style=sameline]
\item[(i)] Place all Physical AI systems in a documented dormancy state, which includes powering down robotic hardware, archiving simulation validation data, and securing all system credentials and access tokens;
\item[(ii)] Maintain Physical AI system maintenance records, calibration data, and software version documentation sufficient to support system reactivation; and
\item[(iii)] Continue cybersecurity monitoring of any Physical AI system components that remain connected to clinical networks during the dormancy period, including patch management for critical security vulnerabilities.
\end{description}

(g) \textit{Physical AI system reactivation.} When a sponsor seeks to resume clinical investigation under an IND that was placed on inactive status and involved Physical AI systems, the protocol amendment required under paragraph (d) shall include, in addition to the standard requirements:

\begin{description}[leftmargin=0.8cm,style=sameline]
\item[(i)] Updated Physical AI System Description reflecting the current hardware configuration, software versions, and AI/ML model versions;
\item[(ii)] Current USL scores for each Physical AI system, reassessed using the most recent version of the USL scoring methodology;
\item[(iii)] Updated simulation validation data demonstrating that the Physical AI system continues to meet cross-framework behavioral consistency requirements after the dormancy period;
\item[(iv)] Updated cybersecurity assessment addressing any vulnerabilities identified during the dormancy period and confirming that all security patches have been applied; and
\item[(v)] Evidence that all investigators and subinvestigators have completed refresher training on the Physical AI system and have demonstrated current competency.
\end{description}

[52 FR 8831, Mar. 19, 1987, as amended at 52 FR 23031, June 17, 1987; 67 FR 9586, Mar. 4, 2002. Physical AI adaptation added 17 March 2026.]

\subsection{\S~312.47 Meetings}

(a) \textit{General.} Meetings between a sponsor and the agency are frequently useful in resolving questions and issues raised during the course of a clinical investigation. FDA encourages such meetings to the extent that they aid in the evaluation of the drug and in the solution of scientific problems concerning the drug, to the extent that FDA's resources permit. The general principle underlying the conduct of such meetings is that there should be free, full, and open communication about any scientific or medical question that may arise during the clinical investigation. These meetings shall be conducted and documented in accordance with part 10.

(b) \textit{``End-of-Phase 2'' meetings and meetings held before submission of a marketing application.}

\begin{description}[leftmargin=0.5cm,style=sameline]
\item[(1)] \textit{End-of-Phase 2 meetings.} The purpose of an end-of-phase 2 meeting is to determine the safety of proceeding to Phase 3, to evaluate the Phase 3 plan and protocols and the adequacy of current studies and plans to assess pediatric safety and effectiveness, and to identify any additional information necessary to support a marketing application. While designed primarily for INDs involving new molecular entities or major new uses, a sponsor of any IND may request and obtain an end-of-Phase 2 meeting. The sponsor should submit background information at least 1 month in advance, including summaries of Phase 1 and 2 investigations, specific Phase 3 protocols, plans for additional nonclinical studies, pediatric study plans, and tentative labeling if available. Agreements reached at the meeting will be recorded in minutes taken by FDA and provided to the sponsor.

\item[(2)] \textit{``Pre-NDA'' and ``pre-BLA'' meetings.} The primary purpose of these meetings is to uncover any major unresolved problems, to identify adequate and well-controlled studies establishing effectiveness, to identify the status of ongoing pediatric studies, to acquaint FDA reviewers with the submission, and to discuss statistical analysis methods and data presentation. The sponsor should submit background information at least 1 month in advance, including a brief summary of clinical studies, a proposed format, information on pediatric studies, and any other information for discussion.
\end{description}

\subsubsection{Physical AI Adaptation of \S~312.47}

(c) \textit{Physical AI system review meetings.} In addition to the meetings described in paragraphs (a) and (b), the following Physical AI-specific meetings are available:

\begin{description}[leftmargin=0.8cm,style=sameline]
\item[(i)] \textit{Pre-IND Physical AI system review.} Prior to submitting an IND that includes a Physical AI system, the sponsor may request a meeting with FDA to review the Physical AI System Description, simulation validation approach, cybersecurity assessment methodology, and proposed human oversight protocols. The sponsor should submit the draft Physical AI System Description at least 1 month in advance.
\item[(ii)] \textit{End-of-Phase 2 Physical AI review.} When an end-of-Phase 2 meeting is requested for an investigation involving Physical AI systems, the sponsor should include in the advance submission materials a summary of Physical AI system performance during Phases 1 and 2, including Physical AI adverse event rates, system reliability metrics, USL score trends, and the proposed Physical AI system configuration for Phase 3 multi-site deployment.
\item[(iii)] \textit{Physical AI system demonstration.} At any meeting described in this section, the sponsor may offer to provide a non-clinical demonstration of the Physical AI system's capabilities, including simulated procedure execution, emergency stop activation, and human-robot interaction workflows. Such demonstrations shall be conducted in a controlled non-clinical environment and shall be documented in the meeting minutes.
\end{description}

[52 FR 8831, Mar. 19, 1987, as amended at 52 FR 23031, June 17, 1987; 55 FR 11580, Mar. 29, 1990; 63 FR 66669, Dec. 2, 1998; 67 FR 9586, Mar. 4, 2002. Physical AI adaptation added 17 March 2026.]

\subsection{\S~312.48 Dispute Resolution}

(a) \textit{General.} The Food and Drug Administration is committed to resolving differences between sponsors and FDA reviewing divisions with respect to requirements for IND's as quickly and amicably as possible through the cooperative exchange of information and views.

(b) \textit{Administrative and procedural issues.} When administrative or procedural disputes arise, the sponsor should first attempt to resolve the matter with the division in FDA's Center for Drug Evaluation and Research or Center for Biologics Evaluation and Research which is responsible for review of the IND, beginning with the consumer safety officer assigned to the application. If the dispute is not resolved, the sponsor may raise the matter with the person designated as ombudsman, whose function shall be to investigate what has happened and to facilitate a timely and equitable resolution.

(c) \textit{Scientific and medical disputes.}

\begin{description}[leftmargin=0.5cm,style=sameline]
\item[(1)] When scientific or medical disputes arise during the drug investigation process, sponsors should discuss the matter directly with the responsible reviewing officials. If necessary, sponsors may request a meeting with the appropriate reviewing officials and management representatives in order to seek a resolution.
\item[(2)] The ``end-of-Phase 2'' and ``pre-NDA'' meetings described in \S~312.47(b) will also provide a timely forum for discussing and resolving scientific and medical issues.
\item[(3)] In requesting a meeting designed to resolve a scientific or medical dispute, applicants may suggest that FDA seek the advice of outside experts, in which case FDA may invite to the meeting one or more advisory committee members or other consultants. For major scientific and medical policy issues not resolved by informal meetings, FDA may refer the matter to a standing advisory committee.
\end{description}

\subsubsection{Physical AI Adaptation of \S~312.48}

(d) \textit{Physical AI technical disputes.} When disputes arise regarding Physical AI system requirements, including USL threshold interpretations, simulation validation adequacy, cybersecurity assessment sufficiency, or the appropriateness of human oversight protocols, the following provisions apply in addition to paragraphs (a) through (c):

\begin{description}[leftmargin=0.8cm,style=sameline]
\item[(i)] \textit{Technical review panel.} For disputes involving Physical AI system safety or performance, the sponsor may request that FDA convene a technical review panel that includes members with expertise in robotic systems engineering, AI/ML validation, simulation physics, cybersecurity, and clinical oncology, in addition to the standard reviewing officials;
\item[(ii)] \textit{Independent simulation validation.} When a dispute centers on the adequacy of simulation validation data, either party may propose that an independent simulation validation study be conducted using the physical-ai-oncology-trials validation benchmark suite (v1.0) or an equivalent validated benchmark, with the results to be reviewed by both parties; and
\item[(iii)] \textit{USL score disputes.} When a dispute centers on a Physical AI system's USL score or the applicable USL threshold for a specific clinical procedure, the sponsor may submit an independent USL assessment conducted by a qualified third party, which FDA shall consider in its review.
\end{description}

[52 FR 8831, Mar. 19, 1987, as amended at 55 FR 11580, Mar. 29, 1990. Physical AI adaptation added 17 March 2026.]

\clearpage

\section{SUBPART D --- RESPONSIBILITIES OF SPONSORS AND INVESTIGATORS}

\subsection{\S~312.50 General Responsibilities of Sponsors}

A sponsor is responsible for selecting qualified investigators, providing them with the information they need to conduct an investigation properly, ensuring proper monitoring of the investigation(s), ensuring that the investigation(s) is conducted in accordance with the general investigational plan and protocols contained in the IND, maintaining an effective IND with respect to the investigations, and ensuring that FDA and all participating investigators are promptly informed of significant new adverse effects or risks with respect to the drug. Additional specific responsibilities of sponsors are described elsewhere in this part.

\subsubsection{Physical AI Adaptation of \S~312.50}

In addition to the general responsibilities described above, a sponsor of a clinical investigation that involves Physical AI systems shall be responsible for:

\begin{description}[leftmargin=0.8cm,style=sameline]
\item[(a)] \textit{Physical AI system selection and qualification.} Selecting Physical AI systems that are appropriate for the intended clinical procedures and that meet or exceed the minimum USL thresholds specified in \S~312.1(e). The sponsor shall maintain documentation of the selection rationale, including comparative analysis of available Physical AI platforms and justification for the selected system's suitability for oncology drug administration;
\item[(b)] \textit{Investigator Physical AI training.} Ensuring that all investigators and subinvestigators who will operate or supervise Physical AI systems have completed the training specified in the Physical AI System Description and have demonstrated competency through documented assessments;
\item[(c)] \textit{Physical AI system deployment and maintenance.} Ensuring that Physical AI systems are properly deployed, configured, calibrated, and maintained at each investigational site in accordance with the manufacturer's specifications and the Physical AI System Description;
\item[(d)] \textit{Physical AI monitoring.} Ensuring that the Physical AI system performance is monitored continuously during clinical procedures and that monitoring data, including telemetry, safety events, and performance metrics, are captured and retained in accordance with the data management plan;
\item[(e)] \textit{Physical AI adverse event reporting.} Ensuring that Physical AI adverse events are identified, documented, and reported to FDA and all participating investigators in accordance with the Physical AI adverse event reporting requirements of \S~312.32(g);
\item[(f)] \textit{Cybersecurity management.} Maintaining the cybersecurity posture of all Physical AI systems throughout the investigation, including timely application of security patches, monitoring for unauthorized access, and incident response in accordance with the cybersecurity assessment submitted under \S~312.23(g)(iv); and
\item[(g)] \textit{Physical AI system updates.} Managing all software updates, AI/ML model updates, and hardware modifications to the Physical AI system in accordance with the PCCP and the amendment requirements of \S~312.30.
\end{description}

[52 FR 8831, Mar. 19, 1987. Physical AI adaptation added 17 March 2026.]

\subsection{\S~312.52 Transfer of Obligations to a Contract Research Organization}

(a) A sponsor may transfer responsibility for any or all of the obligations set forth in this part to a contract research organization. Any such transfer shall be described in writing. If not all obligations are transferred, the writing is required to describe each obligation being assumed by the contract research organization.

(b) A sponsor may not transfer to a contract research organization the obligation to submit an IND.

(c) A contract research organization that assumes any obligation of a sponsor shall comply with the specific regulatory requirements in this chapter applicable to this obligation and shall be subject to the same regulatory action as a sponsor for failure to comply with any obligation assumed under these regulations.

\subsubsection{Physical AI Adaptation of \S~312.52}

(d) \textit{Transfer of Physical AI obligations.} When a sponsor transfers responsibilities involving Physical AI systems to a contract research organization (CRO), the following additional requirements apply:

\begin{description}[leftmargin=0.8cm,style=sameline]
\item[(i)] The written transfer agreement shall explicitly identify which Physical AI-specific obligations are being transferred, including Physical AI system deployment and maintenance, Physical AI adverse event monitoring and reporting, cybersecurity management, simulation validation oversight, and Physical AI system performance monitoring;
\item[(ii)] The CRO shall demonstrate that it possesses or has contracted for the technical expertise necessary to fulfill the transferred Physical AI obligations, including personnel qualified in robotic systems engineering, AI/ML operations, and cybersecurity;
\item[(iii)] The sponsor retains responsibility for submitting all Physical AI-related IND amendments and for all Physical AI adverse event reports to FDA, even when day-to-day Physical AI system management has been transferred to the CRO, unless the transfer agreement explicitly provides otherwise and FDA has been notified; and
\item[(iv)] The transfer agreement shall specify procedures for the CRO to promptly notify the sponsor of any Physical AI safety events, cybersecurity incidents, or USL score changes that may require IND amendment or reporting to FDA.
\end{description}

[52 FR 8831, Mar. 19, 1987, as amended at 67 FR 9586, Mar. 4, 2002. Physical AI adaptation added 17 March 2026.]

\subsection{\S~312.53 Selecting Investigators and Monitors}

(a) \textit{Selecting investigators.} A sponsor shall select only investigators qualified by training and experience as appropriate experts to investigate the drug.

(b) \textit{Control of the drug.} A sponsor shall ship investigational new drugs only to investigators participating in the investigation.

(c) \textit{Providing information to the investigator.} Before the investigation begins, a sponsor (other than a sponsor-investigator) shall provide to each investigator a copy of the investigator's brochure containing the information described in \S~312.23(a)(5) and shall provide the investigator with the information needed to conduct the investigation properly.

(d) \textit{Selecting monitors.} A sponsor shall select a monitor qualified by training and experience to monitor the progress of the investigation.

\subsubsection{Physical AI Adaptation of \S~312.53}

(e) \textit{Physical AI investigator qualifications.} In addition to the requirements of paragraph (a), when selecting investigators for a clinical investigation involving Physical AI systems, the sponsor shall ensure that each investigator:

\begin{description}[leftmargin=0.8cm,style=sameline]
\item[(i)] Has completed the Physical AI system training program specified in the Physical AI System Description and has demonstrated competency through documented assessments, including hands-on operation of the Physical AI system in a non-clinical setting;
\item[(ii)] Has experience or training in supervising robotic-assisted medical procedures, or has completed a supervised training period using the specific Physical AI system under the guidance of an experienced operator; and
\item[(iii)] Has designated at least one qualified backup operator at the investigational site who has also completed the training requirements, to ensure continuous availability of qualified supervision during Physical AI-assisted procedures.
\end{description}

(f) \textit{Physical AI monitor qualifications.} In addition to the requirements of paragraph (d), monitors for investigations involving Physical AI systems shall have training or experience sufficient to evaluate Physical AI system performance data, interpret Physical AI adverse event reports, and assess compliance with Physical AI-specific protocol requirements.

(g) \textit{Physical AI information for investigators.} In addition to the investigator's brochure required under paragraph (c), the sponsor shall provide each investigator with:

\begin{description}[leftmargin=0.8cm,style=sameline]
\item[(i)] The Physical AI System Description or a summary thereof, including system capabilities, limitations, contraindications, and known failure modes;
\item[(ii)] Physical AI system operating manuals and quick reference guides for all procedure types authorized under the protocol;
\item[(iii)] Emergency procedures for Physical AI system malfunctions, including manual override protocols and emergency stop activation procedures; and
\item[(iv)] The Physical AI adverse event reporting requirements, including the classification of Physical AI adverse events and the timelines for reporting.
\end{description}

[52 FR 8831, Mar. 19, 1987, as amended at 52 FR 23031, June 17, 1987; 61 FR 51530, Oct. 2, 1996; 67 FR 9586, Mar. 4, 2002. Physical AI adaptation added 17 March 2026.]

\subsection{\S~312.54 Emergency Research Under \S~50.24 of This Chapter}

(a) The sponsor shall monitor the progress of all investigations involving an exception from informed consent under \S~50.24 of this chapter. When the sponsor receives from the IRB information concerning the public disclosures required by \S~50.24(a)(7)(ii) and (a)(7)(iii) of this chapter, the sponsor promptly shall submit to the IND file and to Docket Number 95S--0158 in the Division of Dockets Management (HFA--305), Food and Drug Administration, 5630 Fishers Lane, rm.~1061, Rockville, MD 20852, copies of the information that was disclosed, identified by the IND number.

(b) The sponsor also shall monitor such investigations to identify when an IRB determines that it cannot approve the research because it does not meet the criteria in the exception in \S~50.24(a) of this chapter or because of other relevant ethical concerns. The sponsor promptly shall provide this information in writing to FDA, all participating investigators, and other sponsors of investigations involving waiver of informed consent for the same drug.

\subsubsection{Physical AI Adaptation of \S~312.54}

(c) \textit{Physical AI in emergency research.} When a clinical investigation involving a Physical AI system is conducted under an exception from informed consent under \S~50.24 of this chapter, the following additional requirements apply:

\begin{description}[leftmargin=0.8cm,style=sameline]
\item[(i)] The public disclosures required under \S~50.24(a)(7)(ii) and (a)(7)(iii) shall include a description of the Physical AI system's role in the investigation, including the types of procedures the Physical AI system will perform and the human oversight protocols in place;
\item[(ii)] The community consultation process shall specifically address community concerns regarding the use of Physical AI systems in emergency medical procedures without prior informed consent; and
\item[(iii)] The protocol shall include heightened Physical AI safety monitoring requirements, including real-time human oversight with immediate manual override capability for all Physical AI-assisted procedures conducted under the emergency research exception.
\end{description}

[62 FR 32479, June 16, 1997, as amended at 67 FR 9586, Mar. 4, 2002. Physical AI adaptation added 17 March 2026.]

\subsection{\S~312.55 Informing Investigators}

(a) Before the investigation begins, a sponsor (other than a sponsor-investigator) shall provide to each investigator a copy of the investigator's brochure containing the information described in \S~312.23(a)(5).

(b) The sponsor shall, as the overall investigation proceeds, keep each participating investigator informed of new observations discovered by or reported to the sponsor on the drug, particularly with respect to adverse effects and safe use. Such information may be distributed to investigators by means of periodically revised investigator brochures, reprints or published studies, reports or letters to clinical investigators, or other appropriate means.

\subsubsection{Physical AI Adaptation of \S~312.55}

(c) \textit{Physical AI system updates to investigators.} In addition to the requirements of paragraphs (a) and (b), the sponsor shall keep each participating investigator informed of:

\begin{description}[leftmargin=0.8cm,style=sameline]
\item[(i)] All Physical AI adverse events reported at other investigational sites, including the nature of the event, the Physical AI system model and software version involved, and any corrective actions taken;
\item[(ii)] Physical AI system software updates, AI/ML model updates, and hardware modifications that may affect the system's performance or safety profile, including the rationale for the update and any re-training or re-validation requirements;
\item[(iii)] Changes in USL scores for the Physical AI system model being used in the investigation, including the reasons for any score changes and the potential impact on clinical procedures;
\item[(iv)] Cybersecurity alerts, vulnerabilities, or incidents affecting the Physical AI system platform, including recommended mitigations; and
\item[(v)] Updated simulation validation data or changes to the sim-to-real gap tolerances that may affect clinical procedure planning.
\end{description}

[52 FR 8831, Mar. 19, 1987. Physical AI adaptation added 17 March 2026.]

\subsection{\S~312.56 Review of Ongoing Investigations}

(a) A sponsor who discovers that an investigator is not complying with the signed agreement (Form FDA 1572), the general investigational plan, or the requirements of this part or other applicable parts shall promptly either secure compliance, or end the investigator's participation in the investigation.

(b) A sponsor shall review and evaluate the evidence relating to the safety and effectiveness of the drug as it is obtained from the investigator. The sponsor shall make such reports to FDA regarding information relevant to the safety of the drug as are required under \S~312.32. The sponsor shall make annual reports on the progress of the investigation in accordance with \S~312.33.

(c) A sponsor shall promptly investigate and report to FDA and to all participating investigators any adverse experience associated with use of the drug that is both serious and unexpected.

(d) A sponsor of a clinical study who has not submitted a marketing application for the drug under study shall maintain records of the composition, manufacture, and controls of the drug, and make such records available for FDA inspection upon request.

\subsubsection{Physical AI Adaptation of \S~312.56}

(e) \textit{Physical AI compliance review.} In addition to the requirements of paragraph (a), the sponsor shall review investigator compliance with Physical AI-specific requirements, including:

\begin{description}[leftmargin=0.8cm,style=sameline]
\item[(i)] Proper use of the Physical AI system in accordance with the approved protocol and Physical AI System Description;
\item[(ii)] Maintenance of Physical AI system logs, telemetry data, and audit trails as required by \S~312.62;
\item[(iii)] Timely reporting of Physical AI adverse events by investigators to the sponsor;
\item[(iv)] Adherence to the human oversight requirements specified in the protocol, including maintenance of the required operator-to-system ratio; and
\item[(v)] Proper implementation of Physical AI system updates and security patches as directed by the sponsor.
\end{description}

(f) \textit{Physical AI performance review.} The sponsor shall continuously review and evaluate Physical AI system performance data across all investigational sites, including:

\begin{description}[leftmargin=0.8cm,style=sameline]
\item[(i)] Aggregate Physical AI adverse event rates and trends;
\item[(ii)] Cross-site comparison of Physical AI system performance metrics;
\item[(iii)] USL score stability and trends;
\item[(iv)] Simulation-to-reality gap measurements; and
\item[(v)] Cybersecurity incident frequency and severity.
\end{description}

The sponsor shall include a Physical AI system performance summary in each annual report submitted under \S~312.33.

[52 FR 8831, Mar. 19, 1987, as amended at 52 FR 23031, June 17, 1987; 67 FR 9586, Mar. 4, 2002. Physical AI adaptation added 17 March 2026.]

\subsection{\S~312.57 Recordkeeping and Record Retention}

(a) A sponsor shall maintain adequate records showing the receipt, shipment, or other disposition of the investigational drug. These records are required to include, as appropriate, the name of the investigator to whom the drug is shipped, and the date, quantity, and batch or code mark of each such shipment.

(b) A sponsor shall maintain complete and accurate records showing any financial interest in accordance with part 54 of this chapter, as required by \S~312.23(a)(11), and \S~312.33(a)(2).

(c) A sponsor shall retain the records and reports required by this part for 2 years after a marketing application is approved for the drug; or, if an application is not approved for the drug, until 2 years after shipment and delivery of the drug for investigational use is discontinued and FDA has been so notified.

(d) A sponsor shall retain reserve samples of any test article and reference standard for a period of not less than 5 years after the completion or discontinuation of the investigation.

\subsubsection{Physical AI Adaptation of \S~312.57}

(e) \textit{Physical AI system records.} In addition to the records required under paragraphs (a) through (d), the sponsor shall maintain the following Physical AI-specific records:

\begin{description}[leftmargin=0.8cm,style=sameline]
\item[(i)] \textit{Physical AI system deployment records.} Complete records of all Physical AI system deployments, including the system model, serial number, hardware configuration, software version, AI/ML model version, and deployment date for each system deployed to each investigational site;
\item[(ii)] \textit{Physical AI system maintenance records.} Records of all Physical AI system maintenance activities, including calibration, repairs, hardware replacements, software updates, AI/ML model updates, and security patches, with dates, descriptions, and the identity of the personnel performing each activity;
\item[(iii)] \textit{Simulation validation records.} Complete records of all simulation validation activities, including the simulation framework used, validation parameters, test results, sim-to-real gap measurements, and any cross-framework behavioral consistency assessments;
\item[(iv)] \textit{Physical AI telemetry archives.} Archived telemetry data from all Physical AI-assisted clinical procedures, including robotic movement trajectories, force/torque measurements, sensor readings, AI decision logs, and operator interaction records. Telemetry data shall be stored in a format that supports retrospective analysis and shall be hash-chained to prevent tampering;
\item[(v)] \textit{USL assessment records.} Records of all USL assessments and reassessments for each Physical AI system, including the raw scoring data, the assessor's identity and qualifications, and the final USL score;
\item[(vi)] \textit{Cybersecurity records.} Records of all cybersecurity assessments, vulnerability scans, penetration tests, security incidents, and remediation actions related to the Physical AI systems; and
\item[(vii)] \textit{Training records.} Records of all Physical AI system training provided to investigators, subinvestigators, and other site personnel, including the training content, dates, and competency assessment results.
\end{description}

(f) \textit{Physical AI record retention.} Physical AI-specific records described in paragraph (e) shall be retained for the same period specified in paragraph (c). Physical AI telemetry archives under paragraph (e)(iv) may be stored in compressed format after the completion of the investigation, provided that the data can be decompressed and rendered in a human-readable format within a reasonable time upon FDA request.

[52 FR 8831, Mar. 19, 1987, as amended at 63 FR 5252, Feb. 2, 1998; 67 FR 9586, Mar. 4, 2002. Physical AI adaptation added 17 March 2026.]

\subsection{\S~312.58 Inspection of Sponsor's Records and Reports}

(a) FDA inspection. A sponsor shall upon request from any properly authorized officer or employee of the Food and Drug Administration, at reasonable times, permit such officer or employee to have access to and copy and verify any records and reports relating to a clinical investigation conducted under this part. Upon written request by FDA, the sponsor shall submit the records or reports (or copies of them) to FDA. The sponsor shall discontinue shipments of the drug to any investigator who has failed to maintain or make available records or reports of the investigation as required by this part.

(b) Controlled substances. If an investigational new drug is a substance listed in any schedule of the Controlled Substances Act (21 U.S.C. 801; 21 CFR part 1308), records concerning shipment, delivery, receipt, and disposition of the drug, which are required to be maintained under this part or other applicable parts of this chapter shall, upon the request of a properly authorized employee of the Drug Enforcement Administration of the Department of Justice, be made available by the investigator or sponsor to whom the request is made, for inspection and copying.

\subsubsection{Physical AI Adaptation of \S~312.58}

(c) \textit{Physical AI system inspection.} In addition to the records and reports described in paragraph (a), the sponsor shall make available to FDA for inspection:

\begin{description}[leftmargin=0.8cm,style=sameline]
\item[(i)] All Physical AI system records described in \S~312.57(e), including telemetry archives, simulation validation records, and cybersecurity records;
\item[(ii)] Access to the Physical AI system software, including source code for custom-developed components directly related to drug administration safety, AI/ML model architectures and training data documentation, and configuration files. The sponsor may request confidential treatment of proprietary software components under applicable FDA procedures;
\item[(iii)] The ability to observe Physical AI system operations in a non-clinical demonstration setting, including execution of test procedures, emergency stop activation, and fault recovery protocols; and
\item[(iv)] Digital access to hash-chained audit trails and cryptographically signed log files, with documentation of the verification procedures for confirming data integrity.
\end{description}

[52 FR 8831, Mar. 19, 1987, as amended at 67 FR 9586, Mar. 4, 2002. Physical AI adaptation added 17 March 2026.]

\subsection{\S~312.59 Disposition of Unused Supply of Investigational Drug}

The sponsor shall assure the return of all unused supplies of the investigational drug from each individual investigator whose participation in the investigation is discontinued or terminated. The sponsor may authorize alternative disposition of unused supplies of the investigational drug provided this alternative disposition does not expose humans to risks from the drug.

\subsubsection{Physical AI Adaptation of \S~312.59}

When the participation of an investigator is discontinued or terminated and the investigation involved Physical AI systems:

\begin{description}[leftmargin=0.8cm,style=sameline]
\item[(a)] The sponsor shall ensure that all Physical AI systems at the investigator's site are placed in safe standby mode and are decommissioned or retrieved in accordance with the decommissioning plan specified in the Physical AI System Description;
\item[(b)] All Physical AI system data, including telemetry archives, system logs, and configuration data, shall be securely transferred to the sponsor or securely destroyed in accordance with the data management plan;
\item[(c)] All Physical AI system access credentials, authorization tokens, and MCP server configurations at the site shall be revoked; and
\item[(d)] The investigator shall certify in writing that all Physical AI system components have been returned, decommissioned, or disposed of in accordance with the sponsor's instructions, and that no Physical AI system at the site retains access to the investigational drug administration systems.
\end{description}

[52 FR 8831, Mar. 19, 1987, as amended at 67 FR 9586, Mar. 4, 2002. Physical AI adaptation added 17 March 2026.]

\subsection{\S~312.60 General Responsibilities of Investigators}

An investigator is responsible for ensuring that an investigation is conducted according to the signed investigator statement, the investigational plan, and applicable regulations; for protecting the rights, safety, and welfare of subjects under the investigator's care; and for the control of drugs under investigation. An investigator shall, in accordance with the provisions of part 50 of this chapter, obtain the informed consent of each human subject to whom the drug is administered, except as provided in \S~50.23 or \S~50.24. Additional specific responsibilities of clinical investigators are set forth in this part and in parts 50 and 56.

\subsubsection{Physical AI Adaptation of \S~312.60}

In addition to the general responsibilities described above, when an investigation involves Physical AI systems, each investigator shall be responsible for:

\begin{description}[leftmargin=0.8cm,style=sameline]
\item[(a)] \textit{Physical AI system operation.} Ensuring that the Physical AI system is operated in accordance with the approved protocol and the Physical AI System Description, including adherence to the authorized procedure types, drug administration parameters, and safety limits;
\item[(b)] \textit{Human oversight.} Maintaining continuous human oversight of the Physical AI system during all clinical procedures involving investigational drugs, including the ability to immediately assume manual control or activate the emergency stop at any time;
\item[(c)] \textit{Pre-procedure verification.} Completing the pre-procedure safety matrix verification before each Physical AI-assisted clinical procedure, confirming system readiness, patient-specific parameter configuration, drug preparation verification, and environmental safety checks;
\item[(d)] \textit{Physical AI adverse event documentation.} Documenting all Physical AI adverse events that occur during clinical procedures under the investigator's supervision, including the event classification, system state at the time of the event, operator actions taken, and patient outcome;
\item[(e)] \textit{Physical AI system integrity.} Reporting to the sponsor any observed Physical AI system anomalies, degraded performance, or unexpected behaviors, even when such observations do not rise to the level of a reportable Physical AI adverse event;
\item[(f)] \textit{Informed consent for Physical AI.} Ensuring that the informed consent process includes a clear explanation of the Physical AI system's role in the investigation, consistent with the requirements of \S~312.23(a)(1)(iv), including a description of the Physical AI system, its function in the clinical procedure, the human oversight protocols in place, the risks associated with the Physical AI system, and the subject's right to request that the procedure be performed without Physical AI assistance if clinically feasible; and
\item[(g)] \textit{Physical AI training currency.} Maintaining current training and competency on the Physical AI system, including completing any refresher training or requalification assessments required by the sponsor following Physical AI system updates.
\end{description}

[52 FR 8831, Mar. 19, 1987. Physical AI adaptation added 17 March 2026.]

\subsection{\S~312.61 Control of the Investigational Drug}

An investigator shall administer the drug only to subjects under the investigator's personal supervision or under the supervision of a subinvestigator responsible to the investigator. The investigator shall not supply the drug to any person not authorized under this part to receive it.

\subsubsection{Physical AI Adaptation of \S~312.61}

When the investigational drug is administered through or with the assistance of a Physical AI system:

\begin{description}[leftmargin=0.8cm,style=sameline]
\item[(a)] The Physical AI system shall be operated only by the investigator or a qualified subinvestigator who has completed the required Physical AI system training and has been authorized by the sponsor;
\item[(b)] The Physical AI system shall not be configured to administer the investigational drug autonomously without the continuous presence and supervision of a qualified investigator or subinvestigator;
\item[(c)] Access controls on the Physical AI system shall be configured to prevent unauthorized persons from initiating drug administration procedures, modifying drug administration parameters, or overriding safety limits; and
\item[(d)] The investigator shall maintain accountability records for the investigational drug that include documentation of any Physical AI-assisted administration, including the date, time, dose, route, Physical AI system identifier, and the identity of the supervising investigator or subinvestigator.
\end{description}

[52 FR 8831, Mar. 19, 1987. Physical AI adaptation added 17 March 2026.]

\subsection{\S~312.62 Investigator Recordkeeping and Record Retention}

(a) \textit{Disposition of drug.} An investigator is required to maintain adequate records of the disposition of the drug, including dates, quantity, and use by subjects. If the investigation is terminated, suspended, discontinued, or completed, the investigator shall return the unused supplies of the drug to the sponsor, or otherwise provide for disposition of the unused supplies of the drug under \S~312.59.

(b) \textit{Case histories.} An investigator is required to prepare and maintain adequate and accurate case histories that record all observations and other data pertinent to the investigation on each individual administered the investigational drug or employed as a control in the investigation. Case histories include the case report forms and supporting data including, for example, signed and dated consent forms and medical records including, for example, progress notes of the physician, the individual's hospital chart(s), and the nurses' notes.

\subsubsection{Physical AI Adaptation of \S~312.62}

(c) \textit{Physical AI procedure records.} In addition to the records required under paragraphs (a) and (b), the investigator shall maintain the following records for each Physical AI-assisted clinical procedure:

\begin{description}[leftmargin=0.8cm,style=sameline]
\item[(i)] The pre-procedure safety matrix verification record, documenting the completion and results of all pre-procedure checks;
\item[(ii)] The Physical AI system identifier (model, serial number, software version, AI/ML model version) used in the procedure;
\item[(iii)] A procedure event log documenting the Physical AI system's actions during the procedure, including drug administration events, operator interventions, safety limit activations, and any system anomalies;
\item[(iv)] Any Physical AI adverse events that occurred during the procedure, documented in accordance with the Physical AI adverse event classification system specified in \S~312.32(g); and
\item[(v)] Post-procedure system status verification, confirming that the Physical AI system is in a safe state and that all procedure data have been captured and stored.
\end{description}

(d) \textit{Physical AI case history integration.} The case histories required under paragraph (b) shall integrate Physical AI system data with clinical observations, including:

\begin{description}[leftmargin=0.8cm,style=sameline]
\item[(i)] Notation that the procedure was Physical AI-assisted and the specific Physical AI system used;
\item[(ii)] Any deviations from the planned Physical AI procedure, including instances where the investigator assumed manual control, activated the emergency stop, or modified Physical AI system parameters during the procedure; and
\item[(iii)] The investigator's clinical assessment of the Physical AI system's performance during the procedure, including any observations relevant to patient safety.
\end{description}

[52 FR 8831, Mar. 19, 1987, as amended at 52 FR 23031, June 17, 1987; 61 FR 57280, Nov. 5, 1996; 67 FR 9586, Mar. 4, 2002. Physical AI adaptation added 17 March 2026.]

\subsection{\S~312.64 Investigator Reports}

(a) \textit{Progress reports.} The investigator shall furnish all reports to the sponsor of the drug who is responsible for collecting and evaluating the results obtained.

(b) \textit{Safety reports.} An investigator shall promptly report to the sponsor any adverse effect that may reasonably be regarded as caused by, or probably caused by, the drug. If the adverse effect is alarming, the investigator shall report the adverse effect immediately.

(c) \textit{Final report.} An investigator shall provide the sponsor with an adequate report shortly after completion of the investigator's participation in the investigation.

\subsubsection{Physical AI Adaptation of \S~312.64}

(d) \textit{Physical AI progress reports.} In addition to the progress reports required under paragraph (a), the investigator shall furnish to the sponsor periodic reports on Physical AI system performance at the investigator's site, including aggregate procedure counts, Physical AI adverse event summaries, system uptime and reliability metrics, and any training or requalification activities completed during the reporting period.

(e) \textit{Physical AI safety reports.} In addition to the safety reports required under paragraph (b):

\begin{description}[leftmargin=0.8cm,style=sameline]
\item[(i)] The investigator shall promptly report to the sponsor any Physical AI adverse event, as defined in \S~312.32(g), regardless of whether the event resulted in patient harm;
\item[(ii)] Physical AI adverse events that are alarming, including uncontrolled robotic movement, drug administration at incorrect doses by the Physical AI system, or loss of emergency stop functionality, shall be reported immediately to the sponsor; and
\item[(iii)] The investigator shall report to the sponsor any observed Physical AI system behavior that, in the investigator's clinical judgment, could compromise patient safety if the behavior were to recur or escalate, even if the observation does not meet the formal definition of a Physical AI adverse event.
\end{description}

(f) \textit{Physical AI final report.} The final report required under paragraph (c) shall include a summary of the investigator's experience with the Physical AI system, including an assessment of the system's reliability, usability, and clinical utility, and any recommendations for Physical AI system improvements based on the clinical investigation experience.

[52 FR 8831, Mar. 19, 1987, as amended at 52 FR 23031, June 17, 1987; 67 FR 9586, Mar. 4, 2002. Physical AI adaptation added 17 March 2026.]

\subsection{\S~312.66 Assurance of IRB Review}

An investigator shall assure that an IRB that complies with the requirements set forth in part 56 will be responsible for the initial and continuing review and approval of the proposed clinical study. The investigator shall also assure that he or she will promptly report to the IRB all changes in the research activity and all unanticipated problems involving risk to human subjects or others, and that he or she will not make any changes in the research without IRB approval, except where necessary to eliminate apparent immediate hazards to human subjects.

\subsubsection{Physical AI Adaptation of \S~312.66}

When a clinical investigation involves Physical AI systems, the investigator shall additionally assure that:

\begin{description}[leftmargin=0.8cm,style=sameline]
\item[(a)] The IRB has been provided with sufficient information about the Physical AI system to conduct a meaningful review of the risks and benefits, including the Physical AI System Description or a summary thereof, the informed consent language regarding the Physical AI system, and a description of the human oversight protocols;
\item[(b)] All changes to the Physical AI system that may affect the risk profile of the investigation, including software updates, AI/ML model changes, and hardware modifications, are reported to the IRB in accordance with its continuing review requirements;
\item[(c)] All Physical AI adverse events that constitute unanticipated problems involving risk to human subjects are promptly reported to the IRB; and
\item[(d)] The IRB is notified of any Physical AI system-related clinical holds, FDA communications regarding Physical AI system deficiencies, or Physical AI-specific amendments to the protocol.
\end{description}

[52 FR 8831, Mar. 19, 1987, as amended at 67 FR 9586, Mar. 4, 2002. Physical AI adaptation added 17 March 2026.]

\subsection{\S~312.68 Inspection of Investigator's Records and Reports}

An investigator shall upon request from any properly authorized officer or employee of FDA, at reasonable times, permit such officer or employee to have access to, and copy and verify any records or reports made by the investigator pursuant to \S~312.62. The investigator is not required to divulge subject names unless the records of particular individuals require a more detailed study of the cases, or unless there is reason to believe that the records do not represent actual cases studied or results obtained.

\subsubsection{Physical AI Adaptation of \S~312.68}

In addition to the records and reports described above, the investigator shall make available to FDA for inspection all Physical AI-specific records required under \S~312.62(c) and (d), including Physical AI procedure event logs, pre-procedure safety matrix records, Physical AI adverse event documentation, and case histories integrating Physical AI system data. Upon FDA request, the investigator shall also facilitate access to the Physical AI system for the purpose of a non-clinical demonstration of system capabilities, emergency stop functionality, and safety protocols.

[52 FR 8831, Mar. 19, 1987, as amended at 67 FR 9586, Mar. 4, 2002. Physical AI adaptation added 17 March 2026.]

\subsection{\S~312.69 Handling of Controlled Substances}

If the investigational drug is subject to the Controlled Substances Act, the investigator shall take adequate precautions, including storage of the investigational drug in a securely locked, substantially constructed cabinet, or other securely locked, substantially constructed enclosure, access to which is limited, to prevent theft or diversion of the substance into illegal channels of distribution.

\subsubsection{Physical AI Adaptation of \S~312.69}

When a Physical AI system is used to administer an investigational drug that is a controlled substance, the following additional precautions apply:

\begin{description}[leftmargin=0.8cm,style=sameline]
\item[(a)] The Physical AI system's drug storage and handling compartments shall be secured with physical locks and electronic access controls that limit access to authorized investigators and subinvestigators only;
\item[(b)] The Physical AI system shall maintain a tamper-evident, hash-chained log of all controlled substance loading, administration, and disposal events, including the identity of the authorizing investigator, the quantity loaded and administered, and any discrepancies; and
\item[(c)] The Physical AI system's drug administration pathway shall be designed to prevent unauthorized diversion, including physical safeguards against unauthorized removal of the drug from the system and software controls that reconcile drug quantities loaded against quantities administered and returned.
\end{description}

[52 FR 8831, Mar. 19, 1987, as amended at 67 FR 9586, Mar. 4, 2002. Physical AI adaptation added 17 March 2026.]

\subsection{\S~312.70 Disqualification of a Clinical Investigator}

(a) If FDA has information indicating that an investigator (including a sponsor-investigator) has repeatedly or deliberately failed to comply with the requirements of this part, part 50, or part 56, or has submitted to the sponsor or FDA misleading reports, FDA will furnish the investigator written notice of the matter, offer an opportunity for explanation, and offer an opportunity for a regulatory hearing under part 16 on the question of whether the investigator is entitled to receive investigational drugs.

(b) On the basis of the administrative record or a regulatory hearing, if the Commissioner determines that the investigator has repeatedly or deliberately violated the requirements, the Commissioner may prohibit the investigator from receiving investigational drugs. Such notice will include a statement of the basis for such determination and may be revoked if adequate assurances of future compliance are provided.

(c) Each IND and each approved application shall include the name of each investigator who has been determined to be ineligible, together with a brief statement of the basis for such determination. This list shall be updated within 15 working days of such determination.

\subsubsection{Physical AI Adaptation of \S~312.70}

(d) \textit{Physical AI grounds for disqualification.} In addition to the grounds specified in paragraph (a), FDA may furnish written notice to an investigator who has:

\begin{description}[leftmargin=0.8cm,style=sameline]
\item[(i)] Repeatedly or deliberately operated a Physical AI system outside the parameters authorized by the approved protocol and Physical AI System Description;
\item[(ii)] Failed to maintain required human oversight of Physical AI system operations during clinical procedures, including operating the Physical AI system without the required qualified supervision;
\item[(iii)] Failed to report Physical AI adverse events to the sponsor as required by this part;
\item[(iv)] Submitted misleading Physical AI system performance data, training records, or Physical AI adverse event reports to the sponsor or to FDA; or
\item[(v)] Allowed unauthorized personnel to operate or access the Physical AI system in connection with the clinical investigation.
\end{description}

(e) \textit{Physical AI competency suspension.} When an investigator is determined to be ineligible under this section based in whole or in part on Physical AI-related grounds, the determination shall specify whether the investigator is prohibited from:

\begin{description}[leftmargin=0.8cm,style=sameline]
\item[(i)] Receiving investigational drugs generally;
\item[(ii)] Receiving investigational drugs only in connection with Physical AI-assisted investigations; or
\item[(iii)] Operating or supervising any Physical AI system in a clinical investigation, while remaining eligible to conduct non-Physical AI investigations.
\end{description}

[52 FR 8831, Mar. 19, 1987, as amended at 52 FR 23031, June 17, 1987; 67 FR 9586, Mar. 4, 2002. Physical AI adaptation added 17 March 2026.]

